\section{Threats to Validity}
\label{sec:threats}

\subsection{Internal Validity}
\label{sec:threats-internal}

Every \texttt{tabpfn\_*} result in this paper comes from one checkpoint,
verified by SHA-256 hash in Section~\ref{sec:methods-implementation},
using Apple's Metal GPU backend. The whole study, XGBoost included, ran on one machine
with every library forced to a single thread to avoid a documented
XGBoost/PyTorch conflict. A different checkpoint, backend, or thread
count could shift specific numbers, but there is no reason to expect
the broader pattern, that robustness varies by method and direction,
to depend on any of them. Each stress cell was repeated across three seeds, enough for the
significance tests reported in Section~\ref{sec:results-significance},
though a study built specifically to estimate variance precisely
would likely use more.

\subsection{External Validity}
\label{sec:threats-external}

The five datasets audited range from 569 to 4{,}601 rows
(Table~\ref{tab:datasets}), well within TabPFN's intended
small-tabular-data regime but far from the scale where gradient-boosted
methods are typically deployed. Whether the same robustness ranking
holds on larger, higher-dimensional, or more categorical-heavy tabular
data is untested. The two corruption mechanisms audited are both
class-conditional (Section~\ref{sec:methods-corruption}). The flip
probability depends only on the true class, not on the specific
feature values of the row being flipped. Real-world label noise is
often instance-dependent instead, arising from cases that are
genuinely ambiguous or that share features with the other class. This
study says nothing about whether the same methods stay robust under
that kind of noise. The four tested contamination rates stop at 0.30,
and results are never extrapolated beyond it
(Section~\ref{sec:methods-contribution}). This study cannot say
whether the same patterns hold at higher contamination.

\subsection{Construct Validity}
\label{sec:threats-construct}

Correction Gain, Correction Harm Rate, Regret, and the break-even rate
are original to this paper. They are reported here only on balanced
accuracy, not on the other nine metrics tracked in
Section~\ref{sec:methods-metrics}. A correction ranked robust by CHR
on balanced accuracy is not guaranteed to be equally robust by the
same measure on AUPRC or minority recall, a point the calibration
results in Section~\ref{sec:results-reliability} already illustrate
for a different metric pair. The Prior Sensitivity Gap's oracle
variant assumes access to the true, pre-corruption class prior,
something no real deployment has. It is a diagnostic construct, not a
claim about achievable performance, and Section~\ref{sec:results-rq4}
found it is not even reliably better than the observed-prior variant
it is meant to bound.
